\section{图像灰度化与灰度直方图}

\subsection{实验过程}
选用指定图像 \texttt{jinx.jpg}，尺寸为 $432\times768$（高$\times$宽）。
利用OpenCV将彩色图像转为灰度图，保持原始尺寸，随后统计0至255各灰度级的像素数，
绘制整幅图像的灰度直方图。

\begin{minipage}{\linewidth}
  \centering
  \includegraphics[width=.485\linewidth]{jinx.jpg}\hfill
  \includegraphics[width=.485\linewidth]{jinx_gray.png}
  \captionof{figure}{指定彩色图像（左）与灰度图（右）。}
\end{minipage}

\subsection{实验结果与分析}
\ReportFigure[128bp]{histogram.png}{整幅图像的灰度直方图，横轴为灰度级，纵轴为像素数。}

灰度图中人物、武器与背景建筑的轮廓仍然可辨，头发与灯光的颜色则转化为不同的亮度。
直方图主要集中在低灰度区域，平均灰度为55.31，出现次数最多的灰度级为49，
低于128的像素占99.37\%，与画面整体较暗的观察一致。
右侧少量高灰度像素对应灯光和局部高光。灰度化保留了主要轮廓与明暗关系，但原有色彩差异不再直接可见。
